%; whizzy paragraph -pdf xpdf -latex ./whizzypdfptex.sh
%; whizzy-paragraph "^\\\\begin{frame}"
% latex beamer presentation.
% platex, latex-beamer でコンパイルすることを想定。 

%     Tokyo Debian Meeting resources
%     Copyright (C) 2009 Junichi Uekawa

%     This program is free software; you can redistribute it and/or modify
%     it under the terms of the GNU General Public License as published by
%     the Free Software Foundation; either version 2 of the License, or
%     (at your option) any later version.

%     This program is distributed in the hope that it will be useful,
%     but WITHOUT ANY WARRANTY; without even the implied warranty of
%     MERCHANTABILITY or FITNESS FOR A PARTICULAR PURPOSE.  See the
%     GNU General Public License for more details.

%     You should have received a copy of the GNU General Public License
%     along with this program; if not, write to the Free Software
%     Foundation, Inc., 51 Franklin St, Fifth Floor, Boston, MA  02110-1301 USA

\documentclass[cjk,dvipdfmx,12pt]{beamer}
\usetheme{Tokyo}
\usepackage{monthlypresentation}

%  preview (shell-command (concat "evince " (replace-regexp-in-string "tex$" "pdf"(buffer-file-name)) "&"))
%  presentation (shell-command (concat "xpdf -fullscreen " (replace-regexp-in-string "tex$" "pdf"(buffer-file-name)) "&"))
%  presentation (shell-command (concat "evince " (replace-regexp-in-string "tex$" "pdf"(buffer-file-name)) "&"))

%http://www.naney.org/diki/dk/hyperref.html
%日本語EUC系環境の時
\AtBeginDvi{\special{pdf:tounicode EUC-UCS2}}
%シフトJIS系環境の時
%\AtBeginDvi{\special{pdf:tounicode 90ms-RKSJ-UCS2}}

\title{Java ポリシーを読んでみた}
\author{まえだ こうへい mkouhei@debian.or.jp\\IRC nick: mkouhei}
\date{2009年4月18日}
\logo{\includegraphics[width=8cm]{image200607/openlogo-light.eps}}

\begin{document}

\frame{\titlepage{}}

\emtext{読むハメになった経緯}
\section{読むハメになった経緯}
\begin{frame}[containsverbatim]{そもそも別の目的がありました}
\begin{center}\bfseries
友人と別件の活動でのなかで、進捗管理が必要に。

\begin{enumerate}
 \item 某所のPCで使えるのがGanttProjectだけ。
 \item オープンソースだからDebianでも使えるベ。
 \item Debianではパッケージになってない！
 \item 「You, ITPしちゃいなよ。」と天の声。
 \item 初めてITPしてみました。
\end{enumerate}
→ まず、ここで最初の脱線。
\end{center}

\end{frame}

\section{}
\begin{frame}[containsverbatim]{２回目の脱線}
 \begin{center}
  \begin{enumerate}
   \item 「次は、Javaポリシー読んで、勉強会のネタにしちゃいなよ。」
   \item 翻訳して、資料作った。← 今ココ。
  \end{enumerate}


 \textbf{\\問題は山積み。}
  \begin{itemize}
   \item まだITPしただけでdebパッケージ化できていない。
   \item GanttProjectを利用できていない。
   \item 肝心の当初の活動自体の進捗も遅れ気味。
  \end{itemize}
 \end{center}
\end{frame}

\section{Java ポリシーの内容}
\begin{frame}{概要}
 大きくこんな構成になっています。
 \begin{enumerate}
  \item 背景説明
  \item ポリシーの内容
  \item 議論すべき問題
  \item Javaパッケージメンテナ向けのアドバイス
 \end{enumerate}
 対象は、Java仮想マシン、Javaコンパイラ、Javaプログラム、Javaライブラリで
 す。
\end{frame}

\section{背景説明}

\begin{frame}{要点は２つ。}
 \begin{itemize}
  \item Java固有の特徴が多いので、Debianポリシーから分けて、サブポリシー
	とて扱われている。
  \item サブポリシー自体へのコメント、提案などは、java-commonパッケージか
	\footnote{\url{java-common@packages.debian.org}}、Debian Javaメーリング
	リスト\footnote{\url{debian-java@lists.debian.org}}に投稿しる。
 \end{itemize}
\end{frame}

\begin{frame}{ポリシー}
Javaパッケージのポリシーは大きく分けて５つあります。
 \begin{enumerate}
  \item 仮想マシン
  \item コンパイラ
  \item プログラム
  \item ライブラリ
  \item main, contrib, non-free
 \end{enumerate}
\end{frame}

\begin{frame}{Javaポリシー}
 \begin{itemize}
  \item Javaパッケージ自体は、コンパイラ、仮想マシン、ランタイム用の仮想
	パッケージが作られます。
  \item Javaのコードはバイトコードとして、全アーキテクチャ向けに配布しな
	ければなりません。
  \item 開発者向けのライブラリと、エンドユーザ向けのプログラムに分類され
	ます。
 \end{itemize}
\end{frame}

\begin{frame}{仮想マシン}
JVMをパッケージ化する際のポリシーです。パッケージ名、依存関係、コマンド
 名を規程しています。\\
 \begin{itemize}
  \item java-virtual-machineパッケージを提供し、java-commonパッケージに
	依存していること。
  \item ランタイム環境も一緒に配布しても良い。
  \item Java2に準拠したランタイムなら、java2-runtimeを提供。
  \item Sun Javaプログラムと互換性がある場合は、/etc/alternatives/javaという名前で提供。
  \item CLASSPATHも事前に定義しておけ。
  \item 仮想マシンのコードを提供する場合は、hogehoge-devという名前にする。
 \end{itemize}
 \end{frame}

\begin{frame}{コンパイラ}
Javaコンパイラをパッケージかする際のポリシー。\\
 \begin{itemize}
  \item java2-compilerパッケージを提供し、java-commonに依存していること。
  \item Sun JavaのJDKと互換性がある場合は、/etc/alternatives/javacという
	名前にする。
  \item コンパイルに必要なJavaのコアクラスをCLASSPATHに設定する。
 \end{itemize}
\end{frame}

\begin{frame}{プログラム}
Javaプログラムをパッケージかする際のポリシー。
 \begin{itemize}
  \item /usr/bin/以下に配置して、実行可能にしておかなければならない。
  \item プログラムはbinfmt\_miscを使用したJavaクラスかラッパーでもよい。
  \item 特定の環境変数がなくても実行できないとアカン。
  \item 実行ファイルのポリシールールに従わないとダメ。
  \item 補助クラスの場合は、/usr/share/javaディレクトリ以下に配置する。
  \item 仮想マシンとランタイムに依存していないとダメ。
  \item プログラムの名前には規則はない。
 \end{itemize}
\end{frame}

\begin{frame}{binfmt\_miscとは}
 Javaプログラムをネイティブコードのように、直接実行させるための仕組み
 のようです。
 \begin{quote}
  Binfmt\_misc はファイルの先頭の何バイトかをあなたの指定したマジックバイ
  トシーケンスと(特定のビットはマスクして)比較し、バイナリタイプを判定し
  ます。
 \end{quote}\footnote{\url{http://www.linux.or.jp/JF/JFdocs/kernel-docs-2.2/binfmt_misc.txt}}

 \begin{quote}
  OS 内で Javaバイナリの直接サポート(中略)、Linux はサポートをしています！
  下記のことを行った後、他のプログラムと同じように Java アプリケーションと
  Java アプレットを実行できます -
 \end{quote}\footnote{\url{http://www.linux.or.jp/JF/JFdocs/kernel-docs-2.2/java.txt}}
\end{frame}

\begin{frame}{ライブラリ}
 Javaライブラリをパッケージ化する際のポリシー。\\
 開発者向け、ユーザ向けのようには分かれていない。
 \begin{itemize}
  \item libHOGE[version]-javaという形式の名前でないとアカン。
  \item jarアーカイブは/usr/share/java以下に配置しないとダメ。
  \item pacage[-extraname]-fullversion.jarの形式にする。
  \item ランタイムには依存しないとダメだが、仮想マシンには依存してはダメ。
 \end{itemize}
 
\end{frame}

\begin{frame}{main, contrib, non-free}
 パッケージはmain, contrib, non-freeのどれあたるかの判断基準です。
 \begin{itemize}
  \item ソースパッケージをnon-freeなツールでしか正しくコンパイルできなけ
	ればmainにはできない。\\
	パッケージ自体はフリーならcontrib。
  \item フリーなコンパイラは、guavac, gcj, jikesがある。
  \item バイナリパッケージがnon-freeな仮想マシンでしか動かせないと、mainに
	はできない。\\
	パッケージ自体がフリーならcontrib。
 \end{itemize}
\end{frame}

\begin{frame}{議論すべき問題}
概要が列挙されています。詳細がよく分からないものもあります。
 \begin{itemize}
  \item リポジトリの名前と存在自体について。
  \item /usr/share/java以下のシンボリックリンクをスクリプトで置き換える
	べきか。
  \item コアクラスの必要性について。
  \item SunのコミュニティーライセンスはDebianで使えるのか？
 \end{itemize}
など。
\end{frame}

\begin{frame}{Javaパッケージメンテナ向けのアドバイス}
Javaパッケージ作成者向けのアドバイス。\\
でも記載内容は、
 \begin{itemize}
 \item 足りないツールをだれか作ってくれ。
  \item 他の言語向けパッケージ作成では自動で行っている部分は(まだ環境が整っ
	てないから)を手動でやれ。
  \item 「これはアドバイスに過ぎず、ポリシーの一部ではない」\\
	…ひどい。
 \end{itemize}
\end{frame}

\begin{frame}{まとめ}
 \begin{itemize}
  \item 通常のポリシーに比べた場合の、Javaポリシーの特徴はどうなのかはよ
	く分からない。
  \item でも、フリーなコンパイラが3つある。\\
	（だから、３つともでパッケージ作成検証を行わないとダメ)
  \item mainに入れるための制約が明確に規定されている。\\
	（今までのDebian勉強会で聞いた限りでは他と変わらない気がする。)
  \item ファイルの配置場所や、命名規約やら、依存関係が一部Java固有っぽい。
 \end{itemize}
\end{frame}

\begin{frame}{今後やること}
 以上を踏まえて、
 \begin{itemize}
  \item ganttprojectのdebパッケージ化を進める。
  \item 今回訳した文章は、debian-docに投げて査読してもらう。
  \item 翻訳のやり方もちゃんと学ぶ。
  \item 本来の目的をちゃんと遂行する。
  \item ヨメに「またDebian？」とイヤな顔をされないように努力する。
 \end{itemize}

\end{frame}

\end{document}

;;; Local Variables: ***
;;; outline-regexp: "\\([ 	]*\\\\\\(documentstyle\\|documentclass\\|emtext\\|section\\|begin{frame}\\)\\*?[ 	]*[[{]\\|[]+\\)" ***
;;; End: ***
